\documentclass[a4paper, french, 11pt]{article}

\usepackage[utf8]{inputenc}
\usepackage[T1]{fontenc}
\usepackage[left=3cm, right=3cm, top=1.5cm, bottom=3cm]{geometry}
\usepackage{lmodern}
\usepackage{xspace}
\usepackage{varioref}
\usepackage{float}
\usepackage{xcolor}
\usepackage{siunitx}
\usepackage{amsmath}
\usepackage{amsthm}
\usepackage{tikz}
\usepackage{amssymb}
\usepackage{tabularx}
\usepackage{enumitem}
\usepackage{babel}
\usepackage{listings}
\usepackage{csquotes}% Recommended

% LISTING CONFIG
\newtoggle{InString}{}% Keep track of if we are within a string
\togglefalse{InString}% Assume not initally in string
\newcommand*{\ColorIfNotInString}[1]{\iftoggle{InString}{#1}{\color{green!60!black}#1}}%
\newcommand*{\ProcessQuote}[1]{#1\iftoggle{InString}{\global\togglefalse{InString}}{\global\toggletrue{InString}}}%
\lstset{
  literate=
  {0}{{{\ColorIfNotInString{0}}}}1
  {1}{{{\ColorIfNotInString{1}}}}1
  {2}{{{\ColorIfNotInString{2}}}}1
  {3}{{{\ColorIfNotInString{3}}}}1
  {4}{{{\ColorIfNotInString{4}}}}1
  {5}{{{\ColorIfNotInString{5}}}}1
  {6}{{{\ColorIfNotInString{6}}}}1
  {7}{{{\ColorIfNotInString{7}}}}1
  {8}{{{\ColorIfNotInString{8}}}}1
  {9}{{{\ColorIfNotInString{9}}}}1
}
\lstset{
  literate=
  {á}{{\'a}}1 {é}{{\'e}}1 {í}{{\'i}}1 {ó}{{\'o}}1 {ú}{{\'u}}1
  {Á}{{\'A}}1 {É}{{\'E}}1 {Í}{{\'I}}1 {Ó}{{\'O}}1 {Ú}{{\'U}}1
  {à}{{\`a}}1 {è}{{\`e}}1 {ì}{{\`i}}1 {ò}{{\`o}}1 {ù}{{\`u}}1
  {À}{{\`A}}1 {È}{{\'E}}1 {Ì}{{\`I}}1 {Ò}{{\`O}}1 {Ù}{{\`U}}1
  {ä}{{\"a}}1 {ë}{{\"e}}1 {ï}{{\"i}}1 {ö}{{\"o}}1 {ü}{{\"u}}1
  {Ä}{{\"A}}1 {Ë}{{\"E}}1 {Ï}{{\"I}}1 {Ö}{{\"O}}1 {Ü}{{\"U}}1
  {â}{{\^a}}1 {ê}{{\^e}}1 {î}{{\^i}}1 {ô}{{\^o}}1 {û}{{\^u}}1
  {Â}{{\^A}}1 {Ê}{{\^E}}1 {Î}{{\^I}}1 {Ô}{{\^O}}1 {Û}{{\^U}}1
  {œ}{{\oe}}1 {Œ}{{\OE}}1 {æ}{{\ae}}1 {Æ}{{\AE}}1 {ß}{{\ss}}1
  {ű}{{\H{u}}}1 {Ű}{{\H{U}}}1 {ő}{{\H{o}}}1 {Ő}{{\H{O}}}1
  {ç}{{\c c}}1 {Ç}{{\c C}}1 {ø}{{\o}}1 {å}{{\r a}}1 {Å}{{\r A}}1
  {€}{{\EUR}}1 {£}{{\pounds}}1
}
\definecolor{codegray}{rgb}{0.5,0.5,0.5}
\definecolor{backcolour}{rgb}{0.95,0.95,0.92}
\lstdefinestyle{myStyle}{
  backgroundcolor=\color{backcolour},
  commentstyle=\color{blue!50!white},
  keywordstyle=\color{orange},
  numberstyle=\small\selectfont\color{codegray},
  stringstyle=\color{red!80!orange},
  basicstyle=\ttfamily\small\selectfont,
  breakatwhitespace=false,
  breaklines=true,
  captionpos=b,
  keepspaces=true,
  numbers=left,
  numbersep=5pt,
  showspaces=false,
  showstringspaces=false,
  showtabs=false,
  tabsize=4,
  xleftmargin=.025\textwidth,
  xrightmargin=.025\textwidth
}


\newtheoremstyle{break}
{}{}
{}{}
{\bfseries}{}
{\newline}{}
\theoremstyle{break}
\newtheorem{Question}{Question}


\newcommand{\po}{\left(}
\newcommand{\pf}{\right)}
\newcommand{\python}{\textsc{Python}\xspace}

\date{}
\title{Gestion des classes en Python}
\author{Martin Andrieux}

\begin{document}
\maketitle

En Python tout est objet, cela peut rendre difficile la compréhension du mécanisme de création de classes, car elle fait elle-même intervenir des classes.

Nous clarifions dans ce document les principaux mécanismes de l'orienté objet en Python, de la création de classe à l'appel de méthode.

\section{Classe \texttt{type}}
Une classe Python est une instance de la classe \texttt{type}.
Ainsi, les deux déclarations suivantes sont équivalentes:

\begin{lstlisting}[style=myStyle, language=python]
class Foo:
  x = 0

class Foo():
  x = 0

Foo = type("Foo", (), {"x" : 0})
\end{lstlisting}


\section{Attributs}
Chaque objet python possède des attributs.
Par exemple, la classe \texttt{Foo} définie précédemment possède un attribut \texttt{x}, tout comme ses éventuelles instances.


Ils sont récupérés dans la sémantique grâce à la fonction \texttt{\_\_getattribute\_\_}, elle même faisant appel à la fonction \textsf{get\_field} de la sémantique

Notons la différence entre \og attributs\fg{} et \og champs\fg{}:
\begin{itemize}
  \item Les \emph{attributs} sont ce que nous manipulons couramment en Python: \texttt{Foo.x} accède à l'attribut \texttt{x} de l'objet \texttt{Foo}.
    Cet accès peut inspecter les superclasses de \texttt{Foo} pour trouver celle qui possède le \emph{champ} \texttt{x}.
  \item Les champs sont directement disponible dans la classe, et ne nécessitent pas d'inspections des superclasses pour être lu.
    Les champs ne sont pas exposés en Python, mais ils sont utilisés dans la sémantique.
    En particulier, se sont les champs qui sont stockés dans le tas à côté de l'objet (voir \textsf{ObjS} dans la sémantique).
\end{itemize}

Les expressions suivantes sont équivalentes:
\begin{lstlisting}[style=myStyle, language=python]
Foo.x
type(Foo).__getattribute__(Foo, "x")
type.__getattribute__(Foo, "x")
\end{lstlisting}

Nous devons donc être capables de retrouver le type d'un objet, pour appeler la méthode \texttt{\_\_getattribute\_\_}.
Cette opération est décrite dans la sémantique (figure 6.7) et est particulièrement simple pour les instances (le type y est directement stocké).

Notons également que \texttt{type.\_\_getattribute\_\_} ne doit pas appeler \texttt{getattribute}, il s'agit d'un cas particulier défini dans la sémantique.

Certains champs ne sont pas stockés dans \textsf{ObjS}.
Ces champs sont dits \og builtin\fg{} et sont récupérés grâce à la fonction \textsf{get\_builtin\_field}.

\section{Superclasses}
Les superclasses sont stockées sous la forme d'une liste de classes, l'ordre des classes dans cette liste est déterminé par \textsf{compute\_mro} (\emph{mro} pour \og Method Resolution Order\fg{})
Il semble s'agir d'un parcours en profondeur de l'arbre d'héritage.
La recherche d'un attribut se fait en parcourant les superclasses dans l'ordre \emph{mro}.

\section{Méthodes}
Une classe contient un ensemble de déclarations, comme la déclaration de \texttt{x} de la classe \texttt{Foo}.
Il est bien sûr possible de déclarer des fonctions, et de les utiliser comme méthodes dans le code.
Toutes ces déclarations sont stockés dans un champs spécial de l'\textsf{ObjN} \textsf{Class}, il s'agit d'un dictionnaire liant les noms aux adresses.

% NOTE:
% Pourquoi ne pas voir ces dictionnaires comme des scopes ?
% Une classe définirait son propre scope, constitué des champs statiques.
% Les superclasses peuvent ainsi être représentées par des `scope_up` (à condition de transformer scope_up en une liste de scopes)
% Fonctionne aussi sans modifier le scope, en stockant les superclasses dans l'ObjN class

Une méthode est une fonction accompagnée de son premier paramètre.
Une méthode est l'application du premier paramètre d'une fonction.
Dans le cas des fonctions de classes, le paramètre \texttt{self} est passé.
Cela permet de l'appliquer de manière transparente lorsque le reste des arguments est passé à la méthode.
Le type \texttt{method} surcharge donc la méthode \texttt{\_\_call\_\_} de la manière suivante:
\begin{lstlisting}[style=myStyle, language=python]
Method(arg1, fun).__call__(*args):
    return f(arg1, *args)
\end{lstlisting}

Les méthodes sont créées avec la fonction \texttt{\_\_get\_\_} du type \texttt{function}.
Il y a équivalence entre les deux dernières expressions du morceau de code suivant:
\begin{lstlisting}[style=myStyle, language=python]
class Foo:
    def bar:
        pass

a.bar
Foo.bar.__get__(a)
\end{lstlisting}

L'appel à \texttt{\_\_get\_\_} implique une allocation mémoire pour stocker la méthode fraîchement créée.
Il y a donc une allocation par occurrence de \texttt{a.bar} dans le code.


\section{Concrètement, dans la sémantique skel}
Pour une déclaration de classe, nous devons évaluer le corps de la classe en récupérant toutes les déclarations.
Ces déclarations sont stocké dans un dictionnaire liant les identifiants aux adresses.
Cela doit pouvoir se faire aisément en utilisant le système de scopes: nous créons un nouveau scope pour l'évaluation du corps de la fonction et récupérons toutes les entrées qu'il contient.
Nous obtenons alors directement le dictionnaire souhaité.

La création d'instance est simple, puisqu'il s'agit simplement d'une référence à la classe.
Les attributs locaux (propres à une instance) sont dans l'\textsf{ObjS} de l'instance, le dictionnaire est initialement vide, puis peuplé par \texttt{\_\_init\_\_}, et les autres méthodes.

La création de méthodes est simple, il faudra généralisé l'appel de fonction aux classes implémentant la méthode \texttt{\_\_call\_\_}.

Le calcul du \emph{mro} n'est pas une priorité, il suffira de modifier la fonction \texttt{getattribute} pour inspecter les superclasses.


\end{document}
